% --- Archon physics formalization source begin ---
% archon:physics
% archon:covers PhyXMiniProblems/problem_phyx_mini_0943.lean
% archon:source-report reports/phyx_mini/problem_phyx_mini_0943.source.json

\chapter{Physics problem phyx\_mini\_0943}
\label{ch:PhyXMiniProblems_problem_phyx_mini_0943}

\paragraph{Problem source.}
\#\# Physical scenario

The series circuit in figure is similar to some radio tuning circuits. It is connected to a variable-frequency ac source with an rms terminal voltage of 1.0\textbackslash{} \textbackslash{}mathrm\{V\}.

\#\# Question

Determine the the inductive reactance \textbackslash{}( X\_L \textbackslash{}), the capacitive reactance \textbackslash{}( X\_C \textbackslash{}), and the impedance \textbackslash{}( Z \textbackslash{})

\#\# Answer choices

- A: A: 200\textbackslash{} \textbackslash{}Omega\textbackslash{}

- B: B: 2000\textbackslash{} \textbackslash{}Omega\textbackslash{}

- C: C: 1000\textbackslash{} \textbackslash{}Omega\textbackslash{}

- D: D: 2200\textbackslash{} \textbackslash{}Omega\textbackslash{}

\#\# Figure caption (auxiliary; use the image as primary evidence)

The image depicts an electrical circuit schematic with various components and measurement labels. Here's a detailed description:

1. **Components:**
   - **Current Source:** A labeled source indicating a current of 2.0 mA.
   - **Voltage Sources:** Five voltage nodes labeled with values: 1.0 V, 1.0 V, 4.0 V, 4.0 V, and 0 V.
   - **Resistor:** Represented by a zigzag line with a resistance value of \textbackslash{}( R = 500 \textbackslash{}, \textbackslash{}Omega \textbackslash{}).
   - **Inductor:** Denoted by a coil symbol with an inductance of \textbackslash{}( L = 0.40 \textbackslash{}, \textbackslash{}text\{mH\} \textbackslash{}).
   - **Capacitor:** Shown as two parallel lines with a capacitance of \textbackslash{}( C = 100 \textbackslash{}, \textbackslash{}text\{pF\} \textbackslash{}).

2. **Connections:**
   - The components are connected in a loop, with labels a, b, c, and d indicating key points in the circuit.
   - The inductor connects between points b and c.
   - The capacitor connects between points c and d.

3. **Diagram Features:**
   - An AC voltage symbol is included at the top, indicating an alternating current in the circuit.
   - Voltage measurement circles are placed at different nodes to indicate the voltage readings there.

This diagram depicts the relationships between these elements in an electrical circuit context.

\#\# Dataset metadata

Electric Circuits; Physical Model Grounding Reasoning, Multi-Formula Reasoning

\paragraph{Recorded answer/context.}
B: B: 2000\textbackslash{} \textbackslash{}Omega\textbackslash{}

\paragraph{Figure/image path.}
/root/proposal\_for\_physic/science-mango/phyx\_mini\_run/phyx\_data/test\_image/943.png

\paragraph{Formalization target.}
create a compiling Lean file with sorry bodies at `PhyXMiniProblems/problem\_phyx\_mini\_0943.lean`. The Lean declarations must preserve the physical quantities, dimensions or dimensional roles, figure labels, governing-law hypotheses, and final relation expressed by this problem.
Use Mathlib/Physlib names found through LeanExplore where available. If a physics API is missing, introduce faithful local abstractions rather than scalar placeholder aliases.

\begin{theorem}[Physics formalization target]
\label{thm:physics:phyx_mini_0943:target}
\lean{PhyXMiniProblems.ProblemPhyXMini0943.problem_phyx_mini_0943}
\uses{def:physics:phyx-mini-0943:phyxminiproblems-problemphyxmini0943-resistancemagnitudeinohms, def:physics:phyx-mini-0943:phyxminiproblems-problemphyxmini0943-seriesrlcsetup, def:physics:phyx-mini-0943:phyxminiproblems-problemphyxmini0943-matchesidealvariablefrequencyseriesrlcscenario, def:physics:phyx-mini-0943:phyxminiproblems-problemphyxmini0943-matchessuppliedseriesrlcfigure, def:physics:phyx-mini-0943:phyxminiproblems-problemphyxmini0943-hasphysicalseriesrlcparameters, def:physics:phyx-mini-0943:phyxminiproblems-problemphyxmini0943-satisfiesidealseriesrlcsteadystatelaws, def:physics:phyx-mini-0943:phyxminiproblems-problemphyxmini0943-recordeddatasetanswer}
The assigned autoformalize agent should translate this physics problem into Lean declarations in the covered file, with theorem and lemma proof bodies written as `by sorry`.
\end{theorem}
\begin{proof}
This is an autoformalization task, not a proof task. Produce statements that can later be proved by the physics prover without weakening the physical model.
\end{proof}

% --- Archon declaration topology begin ---
\section{Declaration topology}
\begin{definition}[electric Current Dimension]
  \label{def:physics:phyx-mini-0943:phyxminiproblems-problemphyxmini0943-electriccurrentdimension}
  \lean{PhyXMiniProblems.ProblemPhyXMini0943.electricCurrentDimension}
  Electric current has physical dimension charge per time
\end{definition}
\begin{proof}
  This is the declared model definition of electric Current Dimension.
\end{proof}
\begin{definition}[electric Potential Dimension]
  \label{def:physics:phyx-mini-0943:phyxminiproblems-problemphyxmini0943-electricpotentialdimension}
  \lean{PhyXMiniProblems.ProblemPhyXMini0943.electricPotentialDimension}
  Electric potential difference has physical dimension energy per charge
\end{definition}
\begin{proof}
  This is the declared model definition of electric Potential Dimension.
\end{proof}
\begin{definition}[electrical Resistance Dimension]
  \label{def:physics:phyx-mini-0943:phyxminiproblems-problemphyxmini0943-electricalresistancedimension}
  \lean{PhyXMiniProblems.ProblemPhyXMini0943.electricalResistanceDimension}
  \uses{def:physics:phyx-mini-0943:phyxminiproblems-problemphyxmini0943-electriccurrentdimension, def:physics:phyx-mini-0943:phyxminiproblems-problemphyxmini0943-electricpotentialdimension}
  Resistance, reactance, and impedance have dimension voltage per current
\end{definition}
\begin{proof}
  This is the declared model definition of electrical Resistance Dimension.
\end{proof}
\begin{definition}[electrical Inductance Dimension]
  \label{def:physics:phyx-mini-0943:phyxminiproblems-problemphyxmini0943-electricalinductancedimension}
  \lean{PhyXMiniProblems.ProblemPhyXMini0943.electricalInductanceDimension}
  \uses{def:physics:phyx-mini-0943:phyxminiproblems-problemphyxmini0943-electricalresistancedimension}
  Inductance has dimension resistance times time
\end{definition}
\begin{proof}
  This is the declared model definition of electrical Inductance Dimension.
\end{proof}
\begin{definition}[electrical Capacitance Dimension]
  \label{def:physics:phyx-mini-0943:phyxminiproblems-problemphyxmini0943-electricalcapacitancedimension}
  \lean{PhyXMiniProblems.ProblemPhyXMini0943.electricalCapacitanceDimension}
  \uses{def:physics:phyx-mini-0943:phyxminiproblems-problemphyxmini0943-electricpotentialdimension}
  Capacitance has dimension charge per voltage
\end{definition}
\begin{proof}
  This is the declared model definition of electrical Capacitance Dimension.
\end{proof}
\begin{definition}[angular Frequency Dimension]
  \label{def:physics:phyx-mini-0943:phyxminiproblems-problemphyxmini0943-angularfrequencydimension}
  \lean{PhyXMiniProblems.ProblemPhyXMini0943.angularFrequencyDimension}
  Angular frequency has inverse-time dimension; radians are dimensionless
\end{definition}
\begin{proof}
  This is the declared model definition of angular Frequency Dimension.
\end{proof}
\begin{definition}[RMSCurrent Quantity]
  \label{def:physics:phyx-mini-0943:phyxminiproblems-problemphyxmini0943-rmscurrentquantity}
  \lean{PhyXMiniProblems.ProblemPhyXMini0943.RMSCurrentQuantity}
  \uses{def:physics:phyx-mini-0943:phyxminiproblems-problemphyxmini0943-electriccurrentdimension}
  A nonnegative, unit-independent RMS electric current
\end{definition}
\begin{proof}
  This is the declared model definition of RMSCurrent Quantity.
\end{proof}
\begin{definition}[RMSVoltage Quantity]
  \label{def:physics:phyx-mini-0943:phyxminiproblems-problemphyxmini0943-rmsvoltagequantity}
  \lean{PhyXMiniProblems.ProblemPhyXMini0943.RMSVoltageQuantity}
  \uses{def:physics:phyx-mini-0943:phyxminiproblems-problemphyxmini0943-electricpotentialdimension}
  A nonnegative, unit-independent RMS potential difference
\end{definition}
\begin{proof}
  This is the declared model definition of RMSVoltage Quantity.
\end{proof}
\begin{definition}[Resistance Magnitude Quantity]
  \label{def:physics:phyx-mini-0943:phyxminiproblems-problemphyxmini0943-resistancemagnitudequantity}
  \lean{PhyXMiniProblems.ProblemPhyXMini0943.ResistanceMagnitudeQuantity}
  \uses{def:physics:phyx-mini-0943:phyxminiproblems-problemphyxmini0943-electricalresistancedimension}
  A nonnegative physical resistance, reactance, or impedance magnitude
\end{definition}
\begin{proof}
  This is the declared model definition of Resistance Magnitude Quantity.
\end{proof}
\begin{definition}[Inductance Quantity]
  \label{def:physics:phyx-mini-0943:phyxminiproblems-problemphyxmini0943-inductancequantity}
  \lean{PhyXMiniProblems.ProblemPhyXMini0943.InductanceQuantity}
  \uses{def:physics:phyx-mini-0943:phyxminiproblems-problemphyxmini0943-electricalinductancedimension}
  A nonnegative, unit-independent physical inductance
\end{definition}
\begin{proof}
  This is the declared model definition of Inductance Quantity.
\end{proof}
\begin{definition}[Capacitance Quantity]
  \label{def:physics:phyx-mini-0943:phyxminiproblems-problemphyxmini0943-capacitancequantity}
  \lean{PhyXMiniProblems.ProblemPhyXMini0943.CapacitanceQuantity}
  \uses{def:physics:phyx-mini-0943:phyxminiproblems-problemphyxmini0943-electricalcapacitancedimension}
  A nonnegative, unit-independent physical capacitance
\end{definition}
\begin{proof}
  This is the declared model definition of Capacitance Quantity.
\end{proof}
\begin{definition}[Angular Frequency Quantity]
  \label{def:physics:phyx-mini-0943:phyxminiproblems-problemphyxmini0943-angularfrequencyquantity}
  \lean{PhyXMiniProblems.ProblemPhyXMini0943.AngularFrequencyQuantity}
  \uses{def:physics:phyx-mini-0943:phyxminiproblems-problemphyxmini0943-angularfrequencydimension}
  A nonnegative, unit-independent physical angular frequency
\end{definition}
\begin{proof}
  This is the declared model definition of Angular Frequency Quantity.
\end{proof}
\begin{definition}[nonnegative SIReadout]
  \label{def:physics:phyx-mini-0943:phyxminiproblems-problemphyxmini0943-nonnegativesireadout}
  \lean{PhyXMiniProblems.ProblemPhyXMini0943.nonnegativeSIReadout}
  Read any nonnegative dimensionful quantity in coherent SI units
\end{definition}
\begin{proof}
  This is the declared model definition of nonnegative SIReadout.
\end{proof}
\begin{definition}[rms Current In Amperes]
  \label{def:physics:phyx-mini-0943:phyxminiproblems-problemphyxmini0943-rmscurrentinamperes}
  \lean{PhyXMiniProblems.ProblemPhyXMini0943.rmsCurrentInAmperes}
  \uses{def:physics:phyx-mini-0943:phyxminiproblems-problemphyxmini0943-rmscurrentquantity, def:physics:phyx-mini-0943:phyxminiproblems-problemphyxmini0943-nonnegativesireadout}
  Read an RMS current in amperes
\end{definition}
\begin{proof}
  This is the declared model definition of rms Current In Amperes.
\end{proof}
\begin{definition}[rms Current In Milliamperes]
  \label{def:physics:phyx-mini-0943:phyxminiproblems-problemphyxmini0943-rmscurrentinmilliamperes}
  \lean{PhyXMiniProblems.ProblemPhyXMini0943.rmsCurrentInMilliamperes}
  \uses{def:physics:phyx-mini-0943:phyxminiproblems-problemphyxmini0943-rmscurrentquantity, def:physics:phyx-mini-0943:phyxminiproblems-problemphyxmini0943-rmscurrentinamperes}
  Read an RMS current in milliamperes
\end{definition}
\begin{proof}
  This is the declared model definition of rms Current In Milliamperes.
\end{proof}
\begin{definition}[rms Voltage In Volts]
  \label{def:physics:phyx-mini-0943:phyxminiproblems-problemphyxmini0943-rmsvoltageinvolts}
  \lean{PhyXMiniProblems.ProblemPhyXMini0943.rmsVoltageInVolts}
  \uses{def:physics:phyx-mini-0943:phyxminiproblems-problemphyxmini0943-rmsvoltagequantity, def:physics:phyx-mini-0943:phyxminiproblems-problemphyxmini0943-nonnegativesireadout}
  Read an RMS potential difference in volts
\end{definition}
\begin{proof}
  This is the declared model definition of rms Voltage In Volts.
\end{proof}
\begin{definition}[resistance Magnitude In Ohms]
  \label{def:physics:phyx-mini-0943:phyxminiproblems-problemphyxmini0943-resistancemagnitudeinohms}
  \lean{PhyXMiniProblems.ProblemPhyXMini0943.resistanceMagnitudeInOhms}
  \uses{def:physics:phyx-mini-0943:phyxminiproblems-problemphyxmini0943-resistancemagnitudequantity, def:physics:phyx-mini-0943:phyxminiproblems-problemphyxmini0943-nonnegativesireadout}
  Read a resistance, reactance, or impedance magnitude in ohms
\end{definition}
\begin{proof}
  This is the declared model definition of resistance Magnitude In Ohms.
\end{proof}
\begin{definition}[inductance In Henries]
  \label{def:physics:phyx-mini-0943:phyxminiproblems-problemphyxmini0943-inductanceinhenries}
  \lean{PhyXMiniProblems.ProblemPhyXMini0943.inductanceInHenries}
  \uses{def:physics:phyx-mini-0943:phyxminiproblems-problemphyxmini0943-inductancequantity, def:physics:phyx-mini-0943:phyxminiproblems-problemphyxmini0943-nonnegativesireadout}
  Read an inductance in henries
\end{definition}
\begin{proof}
  This is the declared model definition of inductance In Henries.
\end{proof}
\begin{definition}[inductance In Millihenries]
  \label{def:physics:phyx-mini-0943:phyxminiproblems-problemphyxmini0943-inductanceinmillihenries}
  \lean{PhyXMiniProblems.ProblemPhyXMini0943.inductanceInMillihenries}
  \uses{def:physics:phyx-mini-0943:phyxminiproblems-problemphyxmini0943-inductancequantity, def:physics:phyx-mini-0943:phyxminiproblems-problemphyxmini0943-inductanceinhenries}
  Read an inductance in millihenries
\end{definition}
\begin{proof}
  This is the declared model definition of inductance In Millihenries.
\end{proof}
\begin{definition}[capacitance In Farads]
  \label{def:physics:phyx-mini-0943:phyxminiproblems-problemphyxmini0943-capacitanceinfarads}
  \lean{PhyXMiniProblems.ProblemPhyXMini0943.capacitanceInFarads}
  \uses{def:physics:phyx-mini-0943:phyxminiproblems-problemphyxmini0943-capacitancequantity, def:physics:phyx-mini-0943:phyxminiproblems-problemphyxmini0943-nonnegativesireadout}
  Read a capacitance in farads
\end{definition}
\begin{proof}
  This is the declared model definition of capacitance In Farads.
\end{proof}
\begin{definition}[capacitance In Picofarads]
  \label{def:physics:phyx-mini-0943:phyxminiproblems-problemphyxmini0943-capacitanceinpicofarads}
  \lean{PhyXMiniProblems.ProblemPhyXMini0943.capacitanceInPicofarads}
  \uses{def:physics:phyx-mini-0943:phyxminiproblems-problemphyxmini0943-capacitancequantity, def:physics:phyx-mini-0943:phyxminiproblems-problemphyxmini0943-capacitanceinfarads}
  Read a capacitance in picofarads
\end{definition}
\begin{proof}
  This is the declared model definition of capacitance In Picofarads.
\end{proof}
\begin{definition}[angular Frequency In Radians Per Second]
  \label{def:physics:phyx-mini-0943:phyxminiproblems-problemphyxmini0943-angularfrequencyinradianspersecond}
  \lean{PhyXMiniProblems.ProblemPhyXMini0943.angularFrequencyInRadiansPerSecond}
  \uses{def:physics:phyx-mini-0943:phyxminiproblems-problemphyxmini0943-angularfrequencyquantity, def:physics:phyx-mini-0943:phyxminiproblems-problemphyxmini0943-nonnegativesireadout}
  Read angular frequency in radians per second
\end{definition}
\begin{proof}
  This is the declared model definition of angular Frequency In Radians Per Second.
\end{proof}
\begin{definition}[Circuit Node]
  \label{def:physics:phyx-mini-0943:phyxminiproblems-problemphyxmini0943-circuitnode}
  \lean{PhyXMiniProblems.ProblemPhyXMini0943.CircuitNode}
  The four electrically distinct junctions labelled in the raster
\end{definition}
\begin{proof}
  This is the declared model definition of Circuit Node.
\end{proof}
\begin{definition}[Circuit Element]
  \label{def:physics:phyx-mini-0943:phyxminiproblems-problemphyxmini0943-circuitelement}
  \lean{PhyXMiniProblems.ProblemPhyXMini0943.CircuitElement}
  The four elements encountered in one traversal of the closed loop
\end{definition}
\begin{proof}
  This is the declared model definition of Circuit Element.
\end{proof}
\begin{definition}[Component Model]
  \label{def:physics:phyx-mini-0943:phyxminiproblems-problemphyxmini0943-componentmodel}
  \lean{PhyXMiniProblems.ProblemPhyXMini0943.ComponentModel}
  Idealized physical roles assigned to the standard circuit symbols
\end{definition}
\begin{proof}
  This is the declared model definition of Component Model.
\end{proof}
\begin{definition}[Voltage Meter Placement]
  \label{def:physics:phyx-mini-0943:phyxminiproblems-problemphyxmini0943-voltagemeterplacement}
  \lean{PhyXMiniProblems.ProblemPhyXMini0943.VoltageMeterPlacement}
  The five voltmeter placements visible in image 943.png
\end{definition}
\begin{proof}
  This is the declared model definition of Voltage Meter Placement.
\end{proof}
\begin{definition}[Series RLCFigure]
  \label{def:physics:phyx-mini-0943:phyxminiproblems-problemphyxmini0943-seriesrlcfigure}
  \lean{PhyXMiniProblems.ProblemPhyXMini0943.SeriesRLCFigure}
  \uses{def:physics:phyx-mini-0943:phyxminiproblems-problemphyxmini0943-circuitnode, def:physics:phyx-mini-0943:phyxminiproblems-problemphyxmini0943-circuitelement, def:physics:phyx-mini-0943:phyxminiproblems-problemphyxmini0943-voltagemeterplacement}
  Literal topology and scalar annotations transcribed from the raster
\end{definition}
\begin{proof}
  This is the declared model definition of Series RLCFigure.
\end{proof}
\begin{definition}[Variable Frequency ACSource]
  \label{def:physics:phyx-mini-0943:phyxminiproblems-problemphyxmini0943-variablefrequencyacsource}
  \lean{PhyXMiniProblems.ProblemPhyXMini0943.VariableFrequencyACSource}
  \uses{def:physics:phyx-mini-0943:phyxminiproblems-problemphyxmini0943-rmsvoltagequantity, def:physics:phyx-mini-0943:phyxminiproblems-problemphyxmini0943-angularfrequencyquantity}
  The variable-frequency source with its independent terminal observable
\end{definition}
\begin{proof}
  This is the declared model definition of Variable Frequency ACSource.
\end{proof}
\begin{definition}[Series RLCSetup]
  \label{def:physics:phyx-mini-0943:phyxminiproblems-problemphyxmini0943-seriesrlcsetup}
  \lean{PhyXMiniProblems.ProblemPhyXMini0943.SeriesRLCSetup}
  \uses{def:physics:phyx-mini-0943:phyxminiproblems-problemphyxmini0943-rmscurrentquantity, def:physics:phyx-mini-0943:phyxminiproblems-problemphyxmini0943-rmsvoltagequantity, def:physics:phyx-mini-0943:phyxminiproblems-problemphyxmini0943-resistancemagnitudequantity, def:physics:phyx-mini-0943:phyxminiproblems-problemphyxmini0943-inductancequantity, def:physics:phyx-mini-0943:phyxminiproblems-problemphyxmini0943-capacitancequantity, def:physics:phyx-mini-0943:phyxminiproblems-problemphyxmini0943-circuitelement, def:physics:phyx-mini-0943:phyxminiproblems-problemphyxmini0943-componentmodel, def:physics:phyx-mini-0943:phyxminiproblems-problemphyxmini0943-voltagemeterplacement, def:physics:phyx-mini-0943:phyxminiproblems-problemphyxmini0943-seriesrlcfigure, def:physics:phyx-mini-0943:phyxminiproblems-problemphyxmini0943-variablefrequencyacsource}
  Independent physical data for the measured series RLC circuit. In particular, the reactance and impedance fields are observables constrained later by constitutive and circuit laws; they are not definitions made from the answer
\end{definition}
\begin{proof}
  This is the declared model definition of Series RLCSetup.
\end{proof}
\begin{definition}[Matches Ideal Variable Frequency Series RLCScenario]
  \label{def:physics:phyx-mini-0943:phyxminiproblems-problemphyxmini0943-matchesidealvariablefrequencyseriesrlcscenario}
  \lean{PhyXMiniProblems.ProblemPhyXMini0943.MatchesIdealVariableFrequencySeriesRLCScenario}
  \uses{def:physics:phyx-mini-0943:phyxminiproblems-problemphyxmini0943-seriesrlcsetup}
  The standard ideal lumped-component interpretation of the shown symbols
\end{definition}
\begin{proof}
  This is the declared model definition of Matches Ideal Variable Frequency Series RLCScenario.
\end{proof}
\begin{definition}[Matches Supplied Series RLCFigure]
  \label{def:physics:phyx-mini-0943:phyxminiproblems-problemphyxmini0943-matchessuppliedseriesrlcfigure}
  \lean{PhyXMiniProblems.ProblemPhyXMini0943.MatchesSuppliedSeriesRLCFigure}
  \uses{def:physics:phyx-mini-0943:phyxminiproblems-problemphyxmini0943-rmscurrentinmilliamperes, def:physics:phyx-mini-0943:phyxminiproblems-problemphyxmini0943-rmsvoltageinvolts, def:physics:phyx-mini-0943:phyxminiproblems-problemphyxmini0943-resistancemagnitudeinohms, def:physics:phyx-mini-0943:phyxminiproblems-problemphyxmini0943-inductanceinmillihenries, def:physics:phyx-mini-0943:phyxminiproblems-problemphyxmini0943-capacitanceinpicofarads, def:physics:phyx-mini-0943:phyxminiproblems-problemphyxmini0943-seriesrlcsetup}
  Primary-image evidence and calibration. The circle marked 2.0 mA is treated as an in-series ammeter, while the five blue circles are RMS voltmeters across the node pairs indicated by their fields
\end{definition}
\begin{proof}
  This is the declared model definition of Matches Supplied Series RLCFigure.
\end{proof}
\begin{definition}[Has Physical Series RLCParameters]
  \label{def:physics:phyx-mini-0943:phyxminiproblems-problemphyxmini0943-hasphysicalseriesrlcparameters}
  \lean{PhyXMiniProblems.ProblemPhyXMini0943.HasPhysicalSeriesRLCParameters}
  \uses{def:physics:phyx-mini-0943:phyxminiproblems-problemphyxmini0943-rmscurrentinamperes, def:physics:phyx-mini-0943:phyxminiproblems-problemphyxmini0943-resistancemagnitudeinohms, def:physics:phyx-mini-0943:phyxminiproblems-problemphyxmini0943-inductanceinhenries, def:physics:phyx-mini-0943:phyxminiproblems-problemphyxmini0943-capacitanceinfarads, def:physics:phyx-mini-0943:phyxminiproblems-problemphyxmini0943-angularfrequencyinradianspersecond, def:physics:phyx-mini-0943:phyxminiproblems-problemphyxmini0943-seriesrlcsetup}
  Positivity and nondegeneracy of the driven passive circuit
\end{definition}
\begin{proof}
  This is the declared model definition of Has Physical Series RLCParameters.
\end{proof}
\begin{definition}[Satisfies Ideal Series RLCSteady State Laws]
  \label{def:physics:phyx-mini-0943:phyxminiproblems-problemphyxmini0943-satisfiesidealseriesrlcsteadystatelaws}
  \lean{PhyXMiniProblems.ProblemPhyXMini0943.SatisfiesIdealSeriesRLCSteadyStateLaws}
  \uses{def:physics:phyx-mini-0943:phyxminiproblems-problemphyxmini0943-rmscurrentinamperes, def:physics:phyx-mini-0943:phyxminiproblems-problemphyxmini0943-rmsvoltageinvolts, def:physics:phyx-mini-0943:phyxminiproblems-problemphyxmini0943-resistancemagnitudeinohms, def:physics:phyx-mini-0943:phyxminiproblems-problemphyxmini0943-inductanceinhenries, def:physics:phyx-mini-0943:phyxminiproblems-problemphyxmini0943-capacitanceinfarads, def:physics:phyx-mini-0943:phyxminiproblems-problemphyxmini0943-angularfrequencyinradianspersecond, def:physics:phyx-mini-0943:phyxminiproblems-problemphyxmini0943-seriesrlcsetup}
  Standard sinusoidal steady-state laws for an ideal series RLC circuit. These are general relations among independently supplied physical observables and contain none of the requested numerical reactance or impedance values
\end{definition}
\begin{proof}
  This is the declared model definition of Satisfies Ideal Series RLCSteady State Laws.
\end{proof}
\begin{definition}[Answer Choice]
  \label{def:physics:phyx-mini-0943:phyxminiproblems-problemphyxmini0943-answerchoice}
  \lean{PhyXMiniProblems.ProblemPhyXMini0943.AnswerChoice}
  Labels of the four resistance-valued choices printed with the problem
\end{definition}
\begin{proof}
  This is the declared model definition of Answer Choice.
\end{proof}
\begin{definition}[displayed Magnitude In Ohms]
  \label{def:physics:phyx-mini-0943:phyxminiproblems-problemphyxmini0943-answerchoice-displayedmagnitudeinohms}
  \lean{PhyXMiniProblems.ProblemPhyXMini0943.AnswerChoice.displayedMagnitudeInOhms}
  \uses{def:physics:phyx-mini-0943:phyxminiproblems-problemphyxmini0943-answerchoice}
  The resistance or reactance magnitude printed beside each choice, in ohms
\end{definition}
\begin{proof}
  This is the declared model definition of displayed Magnitude In Ohms.
\end{proof}
\begin{definition}[recorded Dataset Answer]
  \label{def:physics:phyx-mini-0943:phyxminiproblems-problemphyxmini0943-recordeddatasetanswer}
  \lean{PhyXMiniProblems.ProblemPhyXMini0943.recordedDatasetAnswer}
  \uses{def:physics:phyx-mini-0943:phyxminiproblems-problemphyxmini0943-answerchoice}
  The answer label recorded in the supplied dataset metadata
\end{definition}
\begin{proof}
  This is the declared model definition of recorded Dataset Answer.
\end{proof}
% --- Archon declaration topology end ---
% --- Archon physics formalization source end ---
